%! TEX root = esai.tex
% the main TeX file which is intended to compile, :VimtexReload after adjustment
% :h vimtex-tex-root
\documentclass[../main.tex]{subfiles}
%----------------------------------------------------------------------------------------
%	TITLE SECTION
%----------------------------------------------------------------------------------------
% MAIN TITLE SECTION
\title{
	\textbf{Esai Komitmen Kontribusi ke Instansi Asal/Negara Pasca Studi}\\} % Title and subtitle
%\date{\textbf{\DTMtoday}}
\date{\today}
\author{M. Uliah Shafar}
%\begin{tabular}{@{}ll@{}}
%	Nama & : Muhammad Uliah Shafar\\
%	NIM & : 21020119420029\\
%\end{tabular}
%}

%----------------------------------------------------------------------------------------
% OTHER TITLE SECTION

%\title{\textbf{Sistem Sarana dan Prasarana Jl. Pinggir Laut} \\ {\Large\itshape Infrastructure of Waterfront Parepare City}} % Title and subtitle

%\author{\textbf{Uliah Shafar} \\ \textit{Universitas Diponegoro}} % Author and institution

%\date{\today} % Date, use \date{} for no date

%----------------------------------------------------------------------------------------

\begin{document}

\maketitle

Saya lahir dan besar di sebuah kota  bernama Parepare.
Kota ini adalah satu kota di Indonesia yang tidak pernah lepas dari perhatian banyak orang.
Orang mengenalnya sebagai kota religius, pariwisata, niaga dan pendidikan.
Tidak sedikit juga orang mengetahui bahwa kota ini merupakan tempat kelahiran eyang Habibie, presiden ketiga Republik Indonesia.
% Bakti eyang Habibie dalam ilmu pengetahuan \textit{(science)} patut dihargai apalagi oleh tempat kelahirannya sendiri.
Kepopuleran ini mendorong kota ini untuk terus maju dan berkembang.
% Kepopuleran tersebut menjadi bekal kota Parepare untuk terus berkembang menjadi kota yang lebih baik.

Perkembangan kota Parepare saat ini menunjukkan peningkatan yang signifikan.
Pada akhir tahun 2022, Parepare meresmikan pembangunan kereta api tahap satu yang menghubungkan Parepare dengan Makassar (Ibukota Sulawesi Selatan).
Selain itu, kota ini berhasil menjadikan stadium BJ. Habibie sebagai markas klub sepak bola Persatuan Sepakbola Makassar (PSM).
Pembangunan ini, tidak diragukan, merupakan landasan kota Parepare sebagai magnet bagi ribuan orang.

%Akhir-akhir ini, Parepare mengalami peningkatan urbanisasi yang signifikan.
%can be made explisit permslhan2 urbanisasi di pare
%Berdasarkan data BPS, Parepare mengalami peningkatan penduduk dari tahun 2019 sebesar 5\% di tahun 2021.

Berdasarkan data Badan Pusat Statistik (BPS) tahun 2021, penduduk Parepare meningkat sebesar 5\% dari tahun 2019.
Data ini menunjukkan terjadinya urbanisasi yang berpotensi berpengaruh negatif terhadap masa depan kota. Salah satu dampak negatif tersebut adalah penataan kota menjadi semrawut.
Oleh sebab itu, maka penting untuk merumuskan upaya-upaya dalam menanggulangi dampak tersebut.



%Peningkatan ini menimbulkan permasalahan urbanisasi. Permasalahan ini tentunya dapat berpengaruh terhadap pembangunan parepare di masa depan.
%Oleh karena itu, kota perlu merespon permasalahan tersebut agar tidak berdampak negatif terhadap perkembangan kota. Dalam rangka merespon hal itu, ada sejumlah solusi yang saya jabarkan dalam esai ini.

%Oleh karena itu, kota memerlukan sejumlah solusi terhadap permasalahan tersebut.


% Dalam bidang arsitektur, upaya awal adalah
% Bidang arsitektur perkotaan menjadi salah satu bahan terdepan dalam merumuskan upaya tersebut. Upaya itu dapat melalui pendidikan, praktisi maupun penentu kebijakan.
% Beberapa upaya dalam bidang perkotaan diantaranya adalah,
% pertama,

Dalam bidang arsitektur, upaya-upaya tersebut adalah pertama, meningkatkan kualitas pendidikan terkait perkotaan. Ini dimulai dari tingkat universitas.
Jika kualitas pendidikan di universitas baik,  kemampuan lulusan arsitektur dapat meningkat terutama dalam mengatasi dampak urbanisasi di daerah mereka sendiri. Ini dapat terwujudkan melalui peningkatan kualitas dosen arsitektur di universitas terutama di daerah.
% Dosen menjadi pilar utama dalam meningkatkan pendidikan tersebut. Oleh karena itu, kompetensi dosen-dosen sangat perlu ditingkatkan.

Kedua, mendesain dengan mempertimbangkan  konsep keberlanjutan yang memfokuskan pada pembangunan kota dan permukiman yang inklusif, aman, tangguh dan berlanjut.
Konsep ini akan memberi kota pondasi yang kuat terhadap permasalahan-permasalahan kota khususnya urbanisasi.
Terdapat sejumlah aspek dalam konsep ini yakni aspek ekonomi, energi, ekologi, keikutsertaan dan kesetaraan. Beberapa kota telah mengaplikasikan aspek-aspek tersebut, salah satunya adalah kota Jepara. Kota ini mengaplikasikan setidaknya tiga aspek yaitu aspek ekonomi, energi dan keikutsertaan.

Adapun implementasi dari ketiga aspek tersebut sebagai berikut: pada aspek ekonomi, kota Jepara membangun tempat berkumpul untuk kegiatan-kegiatan bisnis. Tempat tersebut akan mengizinkan orang untuk bersaing secara sehat, membagi ilmu, dan melaksanakan beragam kegiatan dalam perkotaan sehari-hari.
Pada aspek energi, kota menggunakan panel surya sebagai sumber energi di banyak tempat.
Pada aspek keikutsertaan, kota memberdayakan masyarakat dalam menekan dampak dari Pembangkit Listrik Tenaga Uap (PLTU) TJB di Jepara bagian utara.
Implementasi ini tentunya akan mendukung keberlanjutan tidak hanya kota Jepara,
Implementasi ini tentunya dapat diwujudkan pada kota-kota lainnya, sehingga menciptakan keberlanjutan suatu kota secara keseluruhan.

Saya ingin melanjutkan pendidikan hingga ke jenjang doktor karena jenjang ini menyediakan kesempatan untuk mengasah kepekaan terhadap isu manusia dan lingkungan.
Kepekaan tersebut akan mengizinkan saya untuk berkontribusi dalam menanggulangi permasalahan kota seperti urbanisasi di kemudian hari.
Seperti contoh, lanjut mengajar di instansi asal saya dengan kepekaan terhadap isu-isu sekitar yang lebih dalam.

Pada saat kuliah nanti, saya akan mengambil jalur perkuliahan \textit{(by course)}.
Saya akan memfokuskan beberapa mata kuliah berkaitan dengan desain ruang publik, manajemen pengembangan perkotaan, dan arsitektur berkelanjutan dalam perkuliahan.
Selanjutnya, apabila proposal disertasi saya diterima, saya akan mengembangkan judul disertasi tentang identitas ruang publik tepi laut.
Alasannya adalah identitas lokal mulai tergerus dan berubah.
Serta, karakter-karakter suatu tempat yang standar dan seragam semakin berkembang yang sebagian besar dipengaruhi oleh urbanisasi.

% Proposal judul disertasi saya ini telah mendapatkan respon positif dari dosen arsitektur Universitas Hasanuddin yakni Ibu Dr. Ir. Idawarni J. Asmal, M.T. Dia telah menerbitkan satu buku dan sejumlah artikel tentang sebuah identitas perkotaan. Satu penelitiannya berjudul komponen pembentuk kota.
% Berdasarkan itu, beliau menyarankan untuk melihat sebuah identitas kota dari sudut pandang citra kota, teori dari kevin lynch.

Saya berencana melanjutkan studi di universitas dalam negeri. Alasannya adalah universitas dalam negeri sekarang telah memiliki kualitas yang sangat baik. Seperti contoh, ada tiga program studi arsitektur dari universitas dalam negeri yang masuk dalam deretan peringkat 1000 program studi arsitektur terbaik dunia. Prestasi tersebut tentunya dihasilkan dari proses yang menghabiskan banyak waktu dan tenaga.

Ketika saya studi S2 di Universitas Diponegoro, saya merasakan suasana belajar yang penuh semangat. Pada satu kesempatan, dosen menanyai saya tentang makanan orang bugis. Dia bertanya ``Orang sana (suku bugis) tidak makan ikan kan?", respons saya cepat lalu menjawab ``Makan Prof.". Saat itu, beliau membuat kelas bingung terutama saya, lalu menyanggah ``disana tidak makan ikan, cuma makan \textit{ikang}".
Seketika saya mengangguk dan tertawa kecil diikuti oleh teman kelas lainnya. Selain suasana belajar, saya menilai kualitas ilmu pengetahuan yang sangat baik.
Sebagian besar dosen yang mengajari saya adalah doktor, bahkan tidak sedikit yang profesor. Terkadang ada juga dosen luar negeri yang menyempatkan mengajar dalam rangka program \textit{Visiting Professor}. Pengalaman saya ini merupakan gambaran bahwa universitas dalam negeri terus berkembang menjadi lebih baik.

Sekarang, universitas dalam negeri telah melahirkan alumni doktor yang terkemuka dan berkontribusi besar terhadap negeri.
Beberapa diantaranya adalah Moh. Mahfud MD, Agung Gumiwang Kartasasmita, dan Ahmad Basarah. Selain penyandang doktor, banyak tokoh nasional yang telah mengenyam pendidikan di Indonesia, baik sarjana maupun magister. Salah satunya, Bapak Jusuf Kalla, adalah seorang wakil presiden ke-10 dan ke-12. Beliau terkenal di kalangan dunia sebagai tokoh perdamaian. Jusuf kalla, sekarang adalah ketua umum Palang Merah Indonesia (PMI), telah memberikan banyak makna untuk Indonesia.
Tokoh-tokoh ini menandakan bahwa lulusan dalam negeri memiliki nilai dan bekal tersendiri dalam berbakti untuk negeri.

Negara selayaknya kanvas putih yang memerlukan goresan kuas yang beraneka warna. Goresan kuas itu adalah kontribusi anak bangsa yang mulai ada sejak jaman hindia belanda hingga jaman \textit{``now"}. Dulu, umumnya kontribusi tersebut adalah tumpahan darah. Sekarang, kontribusi itu bisa saja sangat beragam. Contohnya, tulisan, desain, aplikasi, suatu gerakan dan lain sebagainya. Goresan pada kanvas tersebut akan membentuk lukisan yang indah dimana itu dinikmati oleh semua orang. % Sebagaiamana peran saya hingga saat ini yang mungkin baru dirasakan oleh segelintir orang.

Sewaktu masih SMA, saya aktif dalam organisasi-organisasi sekolah, khususnya PMR. Pada saat itu, kami mengadakan kegiatan-kegiatan organisasi. Kegiatan yang paling berkesan adalah pertama, saat kami melaksanakan donor darah di sekolah yang diikuti oleh sejumlah guru dan siswa. Kedua, saat kami menyelenggarakan pelatihan-pelatihan untuk pengkaderan seperti pertolongan pertama, penggunaan alat medis umum dan SAR.

Saat kuliah, saya mengikuti organisasi mahasiswa arsitektur di luar kampus yaitu PAMIY (Panguyuban Mahasiswa Arsitektur Yogyakarta). Pada saat itu, kami menyelenggarakan sejumlah kegiatan. Pertama, kami mendesain dan membangun perpustakaan kecil untuk musala di desa Parangritis, Yogyakarta. Kedua, kami menyelenggarakan perlombaan olahraga dan seni antara mahasiswa arsitektur se-Yogyakarta. Selain aktif berorganisasi, saya pernah menjadi asisten dosen untuk mata kuliah \textit{Detail Engineering Design} (DED). Saya membantu mahasiswa untuk mengaplikasikan teori yang baru saja dijelaskan dosen. Terkadang ada mahasiswa yang memerlukan bimbingan lebih lanjut, sehingga saya biasanya menyempatkan untuk mengajar mereka lebih intensif agar dapat mengejar teman lainnya.

Pada saat kuliah S2, saya cenderung aktif dalam menulis daripada mengikuti aktivitas ekstrakurikuler. Alasannya adalah menyebarnya virus yang menghalangi aktivitas-aktivitas seperti biasanya. Meskipun demikian, itu memungkinkan saya untuk mendapatkan IPK 4.00 dan predikat \textit{Cum laude}. Selain itu, saya juga dapat mempublikasikan sebuah artikel di jurnal nasional bereputasi. Tulisan-tulisan termasuk tesis saya mencangkup tentang ruang publik yang saya sangat pedulikan karena berhubungan dengan kehidupan manusia termasuk sosial, kesehatan, dan bahkan kesejahteraan.

Di semester 3 perkuliahan magister saya, saya sempat bekerja di suatu proyek renovasi sekolah madrasah selagi menyelesaikan proposal tesis saya.
Renovasi ini menuntut penambahan desain inklusif dan sebuah taman.
Sehingga, saya sebagai bagian dari juru gambar merancang sebuah ramp sebagaimana yang diminta pada saat peninjauan lokasi di awal pertemuan proyek. Selain itu, kami juga menambahkan taman di depan-depan kelas. Tentunya ini akan menghasilkan desain yang universal artinya penggunaannya / pemanfaatannya melingkupi semua orang secara umum, tanpa batasan fisik, usia dan juga gender.

Beberapa bulan setelah tamat S2, saya bekerja sebagai dosen di Universitas Muhammadiyah Parepare sampai sekarang.
Universitas ini baru saja membuka program studi Perancangan Wilayah dan Kota.
Dalam pekerjaan saya sebagai dosen, saya mendapatkan kesempatan untuk mengampu dua mata kuliah.
Selain itu, saya juga aktif mengikuti kegiatan-kegiatan instansi seperti \textit{workshop} penulisan karya ilmiah untuk dosen muda.

Apabila saya telah selesai menempuh jenjang S3,
saya akan meneruskan pekerjaan saya dengan keterampilan dan wawasan yang lebih baik. Satu pekerjaan dosen yang utama adalah melaksanakan Tri Dharma perguruan tinggi. Beberapa diantaranya pertama, mengajar tentang perkotaan pada program sarjana. Kedua, melakukan penelitian-penelitian yang dapat di terbitkan di jurnal nasional atau internasional. Terakhir, melakukan pengabdian masyarakat di sekitar kawasan Ajatappareng (sejumlah daerah di sekitar Parepare) sebagai penerapan hasil pengajaran dan penelitian.

Semua kontribusi saya tidak terlepas dari pembenahan dari aspek-aspek kepribadian yakni kelebihan, kekurangan, dan pengalaman yang saya miliki.
Pertama, saya merasa memiliki kelebihan yaitu etos kerja yang tinggi, semangat tinggi untuk belajar hal yang baru, dan kegigihan yang kuat. Kedua, saya memiliki kelemahan yakni sifat yang terburu-buru, kurang detail dalam mengerjakan sesuatu, dan perfeksionisme.
Ketiga, saya memiliki pengalaman yang membanggakan yakni mendapatkan predikat \textit{Cum laude} pada saat S2, bekerja pada pembangunan sekolah madrasah, dan diangkat sebagai dosen pns Institut Teknologi Bacharuddin Jusuf Habibie (ITH). Terakhir, pengalaman yang kurang membanggakan saya adalah gagal meraih beasiswa LPDP dan BPI sebelumnya.

Selain itu, kontribusi yang akan saya lakukan pada perguruan tinggi setelah tamat S3 terikat pada program merdeka belajar.
Merdeka belajar merupakan program yang menekankan kebebasan kepada mahasiswa.
Mahasiswa dapat mengambil mata kuliah di luar program studi dan kegiatan diluar lingkup kampus.
Sehingga dosen dengan mudah dapat mengembangkan pembelajaran yang berfokus pada hasil akademik daripada aturan-aturan yang mengikat. Pembenahan aspek kepribadian dan pelaksanaan program-program yang adaptif besar kemungkinan akan memaksimalkan kontribusi saya.

Kontribusi tersebut bertujuan untuk meningkatkan pembelajaran agar misi Indonesia Emas dapat terwujud pada tahun 2045. Oleh karena itu, bukan hanya satu individu, kerja sama dari berbagai pihak akan sangat membantu dalam mencapai tujuan tersebut.
Kerja sama itu antara lain, pertama, peran orangtua dalam membentuk karakter anak sejak kecil. Kedua, peran masyarakat dalam membentuk nilai-nilai yang konstruktif. Terakhir, peran pemerintah dalam mengelola dan memfasilitasi pembelajaran.

Pendidikan yang cemerlang akan menciptakan generasi yang kelak membawa masa depan bangsa yang gemilang.
Seperti contoh, pendidikan ilmu arsitektur atau perkotaan yang baik akan mendorong generasi untuk mengembangkan kota yang bersahabat dengan dampak urbanisasi.
Pendidikan itu dapat terlaksana tidak hanya dengan mengubah kebijakan, namun itu diawali dengan perubahan paradigma dan memerlukan proses panjang.







%Tidak dapat dipungkiri, di masa depan akan muncul peluang-peluang baru dan melimpah dalam menciptakan faedah bagi orang-orang. Seperti contoh,
%saya dapat berkarya dengan cara menulis buku bahkan membuat video tentang arsitektur.
%Lebih dari itu, kemampuan desain dan ilmu pengetahuan yang saya dapatkan akan berguna untuk turut serta mendesain sebuah proyek yang dapat bermanfaat bagi masyarakat.


\end{document}
